\section{地位、与 CSS 的关系、状态表与开放问题}\label{sec:yzliu-status}

本章各陈述的标签、瓶颈与误差类型汇总于\cref{tab:yzliu-status}与\cref{tab:yzliu-status-b}。

\begin{remark}[与 CSS、Collins 与 SCET 的关系]\label{rem:yzliu-relation-css}
本章直接沿用而不重新证明的有：以 Collins 软因子表示的 DY/SIDIS 的 TMD 因子化（\cref{thm:tmd-factorization}，[A]）、CS 演化、单举过程中 Glauber 交换与末态相互作用的相消。超出 CSS 的部分是：软物理的欧氏实现；格点可计算对象的新因子化定理；CSS 型论证隐含使用的引理（时序、解析性、普适类、快度因子化）；重—轻核的阈值与重整子结构；可精确求解的模型；公理化工具（坐标空间解析性、管域、复 Lorentz 协变性、集团性质）。SCET 推导\cite{Vladimirov:2020ofp,Ebert:2022fmh,delRio:2023pse}在把格点软因子等同于 $S_I$ 之后与\cref{thm:lamet-quasi-tmd}一致。全局背景：全局 TMD 拟合\cite{Scimemi:2019cmh}需要的大 $b_\perp$ 输入正是格点纲领想要提供的；Wightman 框架下的 $e^+e^-$ 软因子\cite{Collins:2023cuo}仍是开放的。多部分子软因子快度发散的重整化见第~24~章\cref{prop:front-rapidity-renorm}。
\end{remark}



\begin{tcolorbox}[enhanced,breakable,colback=gray!3!white,colframe=gray!45!black,boxrule=0.5pt,arc=2pt,left=4pt,right=4pt,top=2pt,bottom=2pt,title={常见误解},fonttitle=\bfseries\small]
\small
\begin{itemize}[leftmargin=1.2em,itemsep=1pt,topsep=1pt]
\item 「TMD 的普适性可以直接假定。」不对：SIDIS 与 DY 之间的普适性只到 Wilson 线方向为止，必须由算符定义（订书钉方向）与时间反演推出；本章的时序定理与普适类（\cref{thm:yzliu-time-ordering}、\cref{prop:yzliu-universality-classes}）正是这样的推导。
\item 「单个 TMD 或软函数与快度正规子无关。」不对：只有组合后的对象与正规子无关；欧氏“弯折”软函数从两圈起就不对（\cref{thm:yzliu-soft-hqet}的地位段）。
\item 「TMD 的软因子总是相消、恒为 1。」不对：它只在共线单举截面中为 1；在 TMD 因子化中以平方根进入，而 $S_I$ 是可测的非微扰函数（\cref{prop:yzliu-meson-ff}）。
\item 「快度发散是正规子的人为产物。」不对：物理的重夸克形状因子中 $\ln S$ 随 $Y+Y'$ 线性增长，斜率正是 CS 核（\cref{prop:yzliu-rapidity-linear}、\cref{thm:yzliu-soft-hqet}）。
\item 「TMD 软函数本质上是闵氏的，因此不能在格点上计算。」不对：离光锥软函数是快速色源对的形状因子，有欧氏表示（\cref{thm:yzliu-soft-hqet}）；$S_I$ 也可由介子形状因子提取（\cref{prop:yzliu-meson-ff}）。
\item 「软因子的时序总是无关紧要。」不对：类时双重时序与碎裂型软函数属于不同的普适类（\cref{prop:yzliu-universality-classes}）。
\item 「二维模型的结果表明 QCD 的扭度展开收敛。」不对：\cref{thm:yzliu-gn-convergence}是模型结果，只覆盖虚 Ioffe 时间；动量空间的幂次展开甚至可以发散（\cref{prop:yzliu-on-factorial}）。
\end{itemize}
\end{tcolorbox}

\begin{remark}[开放问题]\label{rem:yzliu-open}
(1) [R] 级的准 TMD／TMD 因子化（奇异闵氏积分的多标度分析）；(2) 把\cref{asm:yzliu-wightman-pt}证明为定理；(3) \cref{conj:yzliu-complex-wilson}以及一致谱隙与光锥一致性假设（\cref{asm:yzliu-lightcone-uniform}）；(4) 类时双重时序软函数的普适性；$e^+e^-$ 软因子；(5) $S_I$ 框架中的极化与胶子准 TMD（部分见\cite{Ebert:2022fmh}）；(6) 大 $\beta_0$ 之外的线性重整子归一化；强子阈值猜想；(7) 四维扭度展开是收敛还是渐近（\cref{thm:yzliu-gn-convergence}与\cref{prop:yzliu-on-factorial}）；(8) 超软与“超—超软”模式的全阶控制，以及\cref{prop:yzliu-rapidity-linear}中的红外步骤；(9) 单圈之外 GPD 在 $x=\pm\xi$ 处的行为。
\end{remark}

\begin{table}[!htbp]
\centering
\footnotesize
\renewcommand{\arraystretch}{1.08}
\begin{tabular}{@{}>{\raggedright\arraybackslash}p{2.25cm}>{\raggedright\arraybackslash}p{3.9cm}>{\raggedright\arraybackslash}p{1.75cm}>{\raggedright\arraybackslash}p{3.15cm}>{\raggedright\arraybackslash}p{1.35cm}>{\raggedright\arraybackslash}p{0.95cm}@{}}
\toprule
陈述 & 内容 & 标签 & 瓶颈 & 误差类型 & 发表 \\
\midrule
\cref{lem:yzliu-partial-exp} & 部分指数化 & [A] & 固定顶点的组合证明未写出 & — & 是 \\
\cref{prop:yzliu-pinch-ultrasoft} & 捏点、超软模式、$L\to\infty$ & [A] & 重求和后超软相消 & E-S & 是 \\
\cref{thm:yzliu-soft-hqet} & 软函数 = 色源对形状因子 & [A] & 欧氏化只在单圈核对 & — & 是 \\
\cref{asm:yzliu-heavy-pair} & 静态 $Q\bar Q$ 谱 & [C] & 禁闭扇区 & — & — \\
\cref{cor:yzliu-soft-form-factor} & 非微扰形状因子形式 & [R, $R_{\mathrm{cond}}$] & \cref{asm:yzliu-heavy-pair} & — & 是 \\
\cref{asm:yzliu-wightman-pt} & 维数正规化中的 Wightman 结构 & [A] & 无精确证明 & — & — \\
\cref{thm:yzliu-analyticity} & 复双曲角解析性 & [A] & \cref{asm:yzliu-wightman-pt} & — & 是 \\
\cref{prop:yzliu-rapidity-linear} & 快度发散线性 & [A] & 最高快度群的红外控制 & E-S & 是 \\
\cref{thm:yzliu-intrinsic-soft} & $S_I$：等价、方案无关、RGE & [A] & 同上 & — & 是 \\
\cref{prop:yzliu-sI-twoloop} & $S_I$ 两圈 & [R, $R_{\mathrm{PT}}$] & — & — & 是 \\
\cref{asm:yzliu-lightcone-uniform} & 光锥极限一致性 & [P] & 只在五条线上 & — & — \\
\cref{prop:yzliu-universality-classes} & 普适类 & [P] & \cref{asm:yzliu-lightcone-uniform} & — & 是 \\
\cref{lem:yzliu-matrix-tree} & 矩阵树公式 & 数学 & — & — & 是 \\
\cref{lem:yzliu-final-state} & 末态求和恒等式 & 数学 & — & — & 是 \\
\cref{thm:yzliu-time-ordering} & 时序无关性 & [A] & \cref{asm:yzliu-wightman-pt} & — & 是 \\
\cref{prop:yzliu-dy-sidis-exp} & 指数正规子下 DY = SIDIS & [P] & 启发式论证 & — & 是 \\
\cref{conj:yzliu-complex-wilson} & 复时空 Wilson 圈 & [C] & 非微扰 & — & 是 \\
\bottomrule
\end{tabular}
\caption{第~23~章状态表（一）：欧氏软函数、解析性、快度因子化与时序定理。“误差类型”列为最弱一步的误差类型（只列出幂次计数给出的 E-S）；“发表”指主要来源是否已在期刊正式发表。}
\label{tab:yzliu-status}
\end{table}

\begin{table}[!htbp]
\centering
\footnotesize
\renewcommand{\arraystretch}{1.08}
\begin{tabular}{@{}>{\raggedright\arraybackslash}p{2.25cm}>{\raggedright\arraybackslash}p{3.9cm}>{\raggedright\arraybackslash}p{1.75cm}>{\raggedright\arraybackslash}p{3.15cm}>{\raggedright\arraybackslash}p{1.35cm}>{\raggedright\arraybackslash}p{0.95cm}@{}}
\toprule
陈述 & 内容 & 标签 & 瓶颈 & 误差类型 & 发表 \\
\midrule
\cref{prop:yzliu-qtmd-powercount} & 准 TMD 幂次计数 & [A] & 多标度分析；Glauber & E-S & 是 \\
\cref{prop:yzliu-hard-kernel} & 匹配核（(b)(c) 为 [R, $R_{\mathrm{PT}}$]） & [A] & 继承准 TMD 定理 & — & 是 \\
\cref{prop:yzliu-lfwf} & 准光前波函数 & [A] & 幂次计数 & E-S & 是 \\
\cref{prop:yzliu-meson-ff} & 介子形状因子与 $S_I$ & [A] & 遍举 Glauber & — & 是 \\
\cref{prop:yzliu-dy-quasi} & 以准 TMD 表示 DY & [A] & 两个 [A] 定理 & — & 是 \\
\cref{prop:yzliu-sivers} & Sivers 比 & [A] & 概要 & — & 是 \\
\cref{thm:yzliu-jet-continuation} & 喷注函数解析延拓 & [A] & \cref{asm:yzliu-wightman-pt}（证明二） & — & 是 \\
\cref{thm:yzliu-threshold} & 阈值因子化 & [A] & 非 Collins 式幂次计数 & — & 是 \\
\cref{prop:yzliu-ad-algebra} & 反常维度代数（两圈 [R, $R_{\mathrm{PT}}$]） & [A] & 继承阈值定理 & — & 是 \\
\cref{prop:yzliu-threshold-nnlo} & NNLO 检验 & [P] & 两个拟合常数 & — & 是 \\
\cref{prop:yzliu-linear-renormalon} & 线性重整子 & [P] & 泡链 & — & 是 \\
\cref{conj:yzliu-renormalon-norm} & $N_m$ 的普适性 & [C] & 超出大 $\beta_0$ & — & 是 \\
\cref{conj:yzliu-threshold-hadronic} & 强子阈值渐近 & [C] & 强子层次 & — & 是 \\
\cref{prop:yzliu-gpd-oneloop} & 准 GPD 单圈结构 & [R, $R_{\mathrm{PT}}$] & — & — & 是$^\dagger$ \\
\cref{conj:yzliu-gpd-xi} & GPD 在 $x=\pm\xi$ & [C] & 软夸克区幂次计数 & — & 否 \\
\cref{thm:yzliu-gn-convergence} & GN 扭度展开收敛 & [P] & 模型；$F_2,F_3$ “类似” & — & 是 \\
\cref{prop:yzliu-gn-structure} & GN 展开结构 & [P] & 模型 & — & 是 \\
\cref{prop:yzliu-on-factorial} & $O(N)$ 阶乘增强 & [P] & 模型 & — & 是 \\
\cref{prop:yzliu-nlsm} & 非线性 $\sigma$ 模型 & [P] & 模型 & — & 是 \\
\cref{prop:yzliu-2dqcd} & 二维 QCD 准 PDF & [P] & 模型 & — & 是 \\
\cref{prop:yzliu-decay-extensions} & 标量理论中的指数衰减 & [P] & 标量、微扰 & — & 否 \\
\bottomrule
\end{tabular}
\caption{第~23~章状态表（续）：准 TMD、遍举与 DY、阈值、重整子、GPD、模型与长程渐近。$^\dagger$单圈核取自已发表的文献\cite{Liu:2019urm}，关于其结构的主张见预印本\cite{Su:2026eta}。}
\label{tab:yzliu-status-b}
\end{table}
